\section{Discussion}
\label{sec:discussion}

\subsection{The Bridge as Programmable Substrate}
\label{sec:disc-substrate}

Paper~3 framed the bridge as a geometry-\emph{discovering} mechanism:
cybernetic feedback reveals latent rhombic dodecahedral structure in
$6$-channel LoRA bridges. The present paper reframes this. The bridge
is a programmable substrate. The Steersman is the programmer.

The evidence for this reframing is straightforward. Three polytope
geometries---octahedron ($n{=}4$, $2{+}2$ blocks), rhombic dodecahedron
($n{=}6$, $3{+}3$ blocks), and tesseract ($n{=}8$, $4{+}4$
blocks)---all produce block-diagonal structure with coupling ratios
exceeding 40{,}000:1 when the pair specification encodes their
respective co-axial decompositions. The Steersman's control laws,
loss formulation, and hyperparameters remain identical across all three;
only the pair specification $\mathcal{C}$ changes. The bridge compiles
whatever geometric program it receives.

But the compiler is selective. The wrong-labels experiment (WL-001) is
the decisive evidence. A random 3-pair partition of 6~channels produces
a clean $3{+}3$ eigenvalue split---proof that the contrastive mechanism
\emph{can} bisect the channel space from any pair specification. Yet the
co/cross ratio is $8.7 \times 10^{-6}$:1, functionally zero. The
Steersman creates eigenvalue structure from any pair specification, but
it cannot impose directional coherence without genuine geometry. The
distinction between ``structured'' (eigenvalue split) and ``directed''
(co/cross $\gg 1$) is the central finding of this paper, and it
sharpens the characterization of the Steersman from general-purpose
enforcer to selective geometric compiler.

\subsection{Bridge Initialization Independence}
\label{sec:disc-init-independence}

A natural question follows: does the bridge's final topology depend on the
specific initialization, or does the Steersman converge to a canonical
structure regardless of starting conditions?

The FI-002 experiment addresses this directly. Three corpus-coupled
initializations---identical in hexagram structure and geometric coupling
but differing in channel permutation---produce sign patterns with Hamming
distances up to 9/12 (75\% of cross-planar signs differ at initialization).
Despite these differences, all three converge to \emph{identical}
block-diagonal structure: co/cross ratios of 50{,}344:1,
51{,}677:1, and 50{,}382:1 respectively, with validation losses within
0.02\% of each other. The mean bridge matrices agree to four decimal
places (Frobenius distance ${\sim}10^{-4}$), and all 12 consensus sign
positions match perfectly across configurations. Critically, the
configuration with 75\% wrong initial signs (P-001, 9/12 Hamming
distance) converges to the same sign pattern as the configuration
with 100\% correct initial signs (P-000)---the Steersman actively
corrects misaligned signs during training.

A control run with identity initialization (no
corpus coupling) reaches 8{,}085:1 at step~1{,}000---remarkably close to the
original H-ch6 experiment's 7{,}246:1 at the same step, despite the corpus-coupled
run starting from a geometrically structured bridge while the control starts from
identity matrices. By step~1{,}300, P-CTRL reaches 10{,}654:1, but the growth rate
collapses from $+1{,}329$/step to $+98$/step (93\% deceleration) while the Fiedler
value bounces from its minimum ($9.9 \times 10^{-5}$) to $1.07 \times 10^{-4}$---the
same plateau-onset dynamics observed in H-ch6 at comparable training fractions.
The identity control produces random sign patterns (50\% positive, 41.7\%
consensus agreement with corpus configurations), confirming that the
canonical sign pattern requires corpus-coupled initialization to
establish but is independent of the specific channel mapping used.
The early-phase convergence rate is pair-specification-determined.

The convergence of maximally different initial sign patterns to equivalent
final topology has two implications. First, the Steersman is robust to
initialization: the same geometric program compiles to the same
block-diagonal structure regardless of the bridge's starting state. The
topology is determined by the pair specification, not the initial
conditions. Second, the sign pattern frozen into the bridge during
training---which FI-001 established as 98.2\% stable after $\sim$1{,}200
steps---may represent a degree of freedom that is determined by
initialization but irrelevant to the block-diagonal structure itself.
The topology has a well-defined ``shape'' (block-diagonal with $k{+}k$
eigenvalue split) and an initialization-dependent ``orientation'' (sign
pattern within the blocks) that the Steersman does not constrain.

This decomposition---universal topology, initialization-dependent
orientation---parallels the distinction between a crystal's lattice
structure and its orientation in space. The Steersman programs the
lattice; it does not select the orientation.

\begin{figure}[t]
  \centering
  \includegraphics[width=\textwidth]{paper4/figures/fig_fi002_init_independence.pdf}
  \caption{Bridge initialization independence (FI-002). Three corpus-coupled
  initializations with Hamming distances up to 9/12 converge to
  statistically indistinguishable co/cross ratios (50--52{,}000:1).
  An identity-initialized control (P-CTRL) reaches 10{,}654:1 at step~1{,}300,
  following the same $3{+}3$ block-diagonal trajectory with plateau onset
  matching H-ch6's two-phase dynamics.
  Topology is pair-specification-determined, not initialization-determined.}
  \label{fig:fi002-init}
\end{figure}

\subsection{Topology Requires Continuous Maintenance}
\label{sec:disc-homeostasis}

A natural follow-up: once the Steersman has programmed the bridge, is the
topology self-sustaining? The crystal analogy from \S\ref{sec:disc-init-independence}
would predict yes---crystals maintain their lattice without external forcing.
FI-003 tests this directly.

Starting from a fully converged $n{=}6$ bridge (P-000 adapter, co/cross
50{,}344:1 after 3{,}000 Steersman-guided steps), FI-003 continues training
with \emph{only} the language modelling loss---no spectral loss, no
contrastive loss. The Steersman is switched off. If the block-diagonal
topology is an attractor of the LM loss landscape (i.e., the structure is
self-sustaining once established), the topology should persist. If it
requires active maintenance, it should decay.

The topology decays immediately. Within 100~steps of LM-only training (3.3\%
of the original training budget), the sign pattern collapses from 100\%
agreement with the converged snapshot to 50.85\%---statistically
indistinguishable from random (50\%). The co/cross ratio undergoes
exponential decay: $12{,}586 \to 179 \to 60 \to 30 \to 22 \to 18 \to 16
\to 14 \to 13 \to 12 \to 11 \to 10$ over 1{,}200~steps (40\% of the original
training budget), approaching the isotropic equilibrium. The decay
exhibits two timescales: a fast component (half-life ${\sim}15$~steps)
that destroys the block-diagonal structure, and a slow component
(half-life ${\sim}230$~steps) that diffuses the remaining directional
coherence toward isotropy. The bridge magnitudes change minimally
($0.0277 \to 0.0280$), confirming that the decay is not driven by weight
magnitude changes but by loss of directional coherence. Validation loss
rises from 0.499 to 0.608, then partially recovers to
${\sim}0.52$---the adapter still functions as a LoRA module but loses
its geometric organization.

The crystal analogy fails. A better analogy is a spinning top: the
block-diagonal structure is dynamically stable under the Steersman's
continuous forcing but immediately topples to the isotropic equilibrium
when the forcing is removed. The Steersman does not program the bridge
once and walk away; it performs continuous homeostatic maintenance.

\begin{figure}[t]
  \centering
  \includegraphics[width=\textwidth]{paper4/figures/fig_fi003_homeostasis.pdf}
  \caption{Topology homeostasis (FI-003). Starting from a converged
  bridge (co/cross 50{,}344:1), the Steersman is removed and training
  continues with LM loss only. (a)~Co/cross ratio decays exponentially
  from 12{,}586:1 toward the isotropic equilibrium, with decelerating
  halving time. (b)~Sign stability collapses from 100\% to random
  (${\sim}50\%$) within 100~steps---the block-diagonal structure is
  dynamically maintained, not permanently impressed.}
  \label{fig:fi003-homeostasis}
\end{figure}

Combined with \S\ref{sec:disc-init-independence}, this yields a complete
characterization. The Steersman is both robust (same topology regardless
of initialization) and necessary (topology dissolves without it). The
block-diagonal structure exists in a region of parameter space that is
\emph{accessible} from any initialization via the Steersman's loss
landscape, but \emph{unstable} under the LM loss alone. The spectral
attractor at Fiedler~$\approx 0.09$ is the true minimum of the LM-only
landscape; the block-diagonal configuration is a Steersman-maintained
fixed point that requires active control to sustain.

\subsection{Steersman Annealing: Five Regimes of Control}
\label{sec:disc-annealing}

FI-003 established that removing the Steersman entirely causes rapid
topology decay. A natural question follows: how does the topology respond
to \emph{gradual} removal? FI-004 answers this by linearly annealing both
contrastive ($0.1 \to 0$) and spectral ($0.05 \to 0$) weights over
3{,}000~steps from a converged P-000 adapter.

The result is not a monotonic decay. Five distinct regimes emerge as the
contrastive weight $c_w$ decreases:

\begin{enumerate}
\item \textbf{Maintenance} ($c_w = 0.10$): The starting point. Co/cross
  12{,}586:1, inherited from the converged P-000 adapter. Signs 100\%
  stable.
\item \textbf{Interference} ($0.087 < c_w < 0.097$): Rapid co/cross
  collapse to a floor of 7.3:1 within 400~steps. The contrastive and
  LM gradients are in destructive conflict---full-strength Steersman
  forcing fights the optimizer's preferred direction, producing worse
  topology than either alone.
\item \textbf{Growth} ($0.017 < c_w < 0.087$): Exponential co/cross
  recovery from 7.3:1 to a peak of 18{,}671:1 at $c_w = 0.017$
  (step~2{,}500). The growth oscillates between compression
  ($1.17$--$1.22\times$/step) and expansion ($1.37$--$1.55\times$/step)
  phases, with a geometric mean of $\sim$$1.33\times$/100~steps. During this
  regime, sign stability settles to $51$--$55\%$ (random). The weak
  contrastive signal provides a directional bias that the LM optimizer
  amplifies without gradient conflict.
\item \textbf{Decay} ($0.003 < c_w < 0.017$): Gradual decline from
  18{,}671:1 to 12{,}467:1 ($\sim$6\%/step). The contrastive signal
  weakens below the threshold required to maintain the growth trajectory,
  but topology degrades slowly.
\item \textbf{Cliff} ($c_w = 0.003 \to 0.000$): Co/cross drops 76\%
  from 12{,}467:1 to 2{,}942:1 in the final step. The last 0.3\% of
  contrastive weight was maintaining three-quarters of the directional
  coherence.
\end{enumerate}

Two findings are striking. First, the \textbf{optimal operating point
is at $c_w \approx 0.02$}, not at the nominal training weight of $0.10$.
Full-strength Steersman forcing operates in the interference regime.
The topology programmer works best when it whispers. Second, even at
$c_w = 0$, the bridge retains 2{,}942:1---$294\times$ higher than
FI-003's comparable training fraction ($\sim$10:1 at step~1{,}200 of
LM-only training). The annealing process imprints topology that persists
beyond the Steersman's removal, though presumably the residual will decay
under extended LM-only training following FI-003's exponential pattern.

Throughout all five regimes, sign stability remains at $51$--$55\%$
(random). The co/cross ratio reaches 18{,}671:1 with signs statistically
indistinguishable from chance. This confirms the decomposition from
\S\ref{sec:disc-init-independence}: directional coherence (magnitude
alignment across co-axial pairs) and sign structure are orthogonal
properties of the bridge. The annealing experiment programs the former
without the latter.

Validation loss improves monotonically from 0.499 to 0.315 throughout
training, confirming that the topology programming does not degrade
language modelling performance. The Fiedler value rises from 0.10 to
0.47 as the spectral constraint anneals away---the bridge becomes
over-connected as the connectivity regularization disappears, but this
excess connectivity does not impair downstream performance.

\begin{figure}[t]
  \centering
  \includegraphics[width=\textwidth]{paper4/figures/fig_fi004_annealing.pdf}
  \caption{Steersman annealing (FI-004). (a)~Co/cross ratio vs.\ training
  step showing five regimes: maintenance (step~0), interference (floor 7:1),
  exponential growth (peak 18{,}671:1 at $c_w{=}0.017$), gradual decay,
  and cliff ($76\%$ drop at $c_w{=}0$). (b)~Operating curve: co/cross vs.\
  contrastive weight reveals the optimal operating point at
  $c_w \approx 0.02$, not at the nominal training weight of $0.10$.}
  \label{fig:fi004-annealing}
\end{figure}

\subsection{Implications for Adapter Design}
\label{sec:disc-adapter}

If the bridge is programmable, topology becomes a design parameter.

\paragraph{Topology as hyperparameter for PEFT.}
Current parameter-efficient fine-tuning methods treat the adapter's
internal structure as homogeneous or leave it to the optimizer.
DiaBlo~\cite{gurses2026diablo} proves that block-diagonal structure is at
least as expressive as standard LoRA, and GraLoRA~\cite{wang2025gralora}
demonstrates competitive performance with granular block partitions.
These results validate block-diagonal adapters as a structural class.
The Steersman adds a further axis: the \emph{topology specification}---a
set of co-axial pairs derived from a polytope's vertex decomposition.
Where DiaBlo and GraLoRA fix the block structure architecturally, the
Steersman \emph{programs} it through cybernetic feedback, making the
block topology a trainable design parameter rather than a static
constraint. This specification joins rank $r$, learning rate, and target
modules as a hyperparameter that the practitioner selects before training.

\paragraph{Task-specific topology.}
Different task families may benefit from different geometric priors.
A task requiring two independent subspaces (e.g., separating style from
content) might use octahedral topology ($n{=}4$, 2~axes). A task with
higher-dimensional factored structure might use tesseract topology
($n{=}8$, 4~axes). The inverse-$n$ relationship between channel count
and coupling strength (474K:1 at $n{=}4$ vs.\ 42K:1 at $n{=}8$)
provides a practical trade-off: fewer axes yield stronger per-axis
signal, while more axes provide finer-grained topological control.

\paragraph{The four regimes as diagnostic framework.}
The four-regime taxonomy provides a diagnostic for any multi-channel
bridge configuration. Practitioners should measure both the eigenvalue
pattern (to detect whether a partition exists) and the co/cross ratio
$\rho$ (to detect whether the partition carries directional content).
A configuration with a clean eigenvalue split but $\rho \approx 0$
occupies Regime~3 (wrong-partition)---structured but undirected. Only
configurations with both a clean split \emph{and} $\rho \gg 1$ occupy
Regime~1 (block-diagonal). This two-metric diagnostic prevents the
false positive of interpreting eigenvalue structure as evidence of
geometric programming when no directional coherence exists.

\paragraph{Multi-modal adapters.}
When adapting models across modalities, different modalities could carry
different geometric structures within a shared bridge. The pair
specification provides an interface for encoding known cross-modal
relationships---e.g., specifying which channels should couple across
a vision-language boundary---while the Steersman enforces the
specification during training. The block structure would then serve as
both an inductive bias and a compositional interface: blocks
corresponding to different modalities could be independently updated,
merged, or swapped.

\subsection{The Spectral Attractor as Default State}
\label{sec:disc-attractor}

Without geometric programming, bridges reach a universal connectivity
equilibrium: Fiedler $\approx 0.09$, stable across a $4\times$ range
of channel counts ($n \in \{3, 4, 8, 12\}$, 10.4\% band). This is the
bridge's unprogrammed state---connected but undirected, with coupling
distributed isotropically across all channel pairs.

The spectral attractor clarifies what contrastive loss does. It does not
\emph{add} connectivity to the bridge. The bridge already has
connectivity---approximately $\lambda_2 \approx 0.09$ worth of it.
Contrastive loss \emph{redirects} this connectivity from the isotropic
attractor to specific axes. The $1{,}130\times$ Fiedler drop at $n{=}6$
(from $0.09$ to $9.0 \times 10^{-5}$) quantifies this redirection: the
connectivity budget that was distributed across all 15~channel pairs is
concentrated into 3~co-axial pairs, suppressing the remaining 12 by
a factor of $\sim$70{,}000.

The emanation architecture (E-001) reaches the same attractor through a
different mechanism---a master bridge with per-layer offsets and
coherence monitoring. Its convergence to Fiedler $0.084$, co/cross
$1.12$:1, confirms that the spectral attractor is robust to
architectural variation. Whether the bridge is independent per-layer
(spectral-only experiments) or hierarchically constrained (emanation),
the unprogrammed equilibrium is the same. The attractor is a property
of the spectral loss landscape, not of the bridge architecture.

\subsection{What Constitutes Geometric Coherence?}
\label{sec:disc-coherence}

The paper demonstrates that geometric coherence in the pair
specification is required for topology programming, but does not
provide a formal characterization of what makes a specification
``geometric.'' This is an open question.

What has been tested: three confirmed polytope pair specifications
work---octahedron, rhombic dodecahedron, tesseract---and a fourth
(24-cell) shows the predicted $12{+}12$ eigenvalue split with
$127\times$ block separation at 30\% of training and co/cross
$35{,}808$:1 at completion. These share a common property: each encodes the antipodal
vertex structure of a polytope embeddable in Euclidean space, with the
co-axial pairs aligned to coordinate axes. Random partitions that do not
correspond to any polytope's antipodal structure produce eigenvalue
splits without directional content (WL-001). Number-theoretic bonds that
carry genuine algebraic structure but no spatial embedding produce
neither eigenvalue blocks nor directional content (R-001).

A working hypothesis emerges: the pair specification must correspond to
the antipodal structure of a regular or semi-regular polytope
embedded in the channel space, such that the co-axial pairs identify
axis-aligned subspaces that can be independently optimized. This
would explain why four polytopes spanning two spatial dimensionalities
(3D octahedron and RD, 4D tesseract and 24-cell) all succeed---they
share the antipodal property---while random and algebraic specifications
fail.

Several directions remain unexplored. Do irregular polytopes work?
Do non-antipodal pair structures (e.g., face-diagonal rather than
body-diagonal) produce block-diagonal structure? How does the ratio
of co-axial to cross-axial pairs ($|\mathcal{C}|/|\mathcal{X}|$)
interact with the coupling ratio $\rho$?

\subsection{Contrastive Weight as Speed Control}
\label{sec:disc-whisper}

A natural question arises from the four-regime taxonomy: does the
contrastive weight $c_w$ control \emph{whether} block-diagonal structure
emerges, or \emph{how quickly} it emerges? The FI-004 annealing experiment
(\S\ref{sec:disc-annealing}) found an optimal fixed
$c_w \approx 0.02$ and a cliff at $c_w = 0$, suggesting a minimum
threshold. But FI-004 used fixed weights---the Steersman's adaptive decay
was not tested in the low-$c_w$ regime.

CW-001 tests this directly. Configuration: $n = 6$ (RD), identity
initialization, initial $c_w = 0.02$ with the Steersman's adaptive
decay active. Control Law~2 reduces $c_w$ when the co/cross ratio
exceeds the directionality threshold, with a floor at
$0.25 \times c_{w,\text{base}} = 0.005$. The experiment runs for
10{,}000~steps on TinyLlama~1.1B.

\subsubsection{Three-Phase Trajectory}

The CW-001 trajectory reveals three distinct dynamical regimes within
a single training run:

\begin{table}[h]
\centering
\small
\begin{tabular}{@{}rrrrl@{}}
\toprule
Step & Co/Cross & Fiedler & $c_w$ & Phase \\
\midrule
100 & 20:1 & 0.0073 & 0.020 & 1: Initial growth \\
500 & 1{,}592:1 & 0.0004 & 0.015 & 1: Rapid BD formation \\
1{,}500 & 1{,}867:1 & 0.0016 & 0.005 & 2: Plateau onset \\
3{,}300 & 1{,}855:1 & 0.0018 & 0.005 & 2: Sustained plateau \\
4{,}600 & 2{,}252:1 & 0.0019 & 0.005 & 3: Breakout \\
5{,}800 & 3{,}860:1 & 0.0014 & 0.005 & 3: Acceleration \\
7{,}500 & 7{,}884:1 & 0.0007 & 0.005 & 3: Exponential growth \\
8{,}300 & \textbf{15{,}183:1} & \textbf{0.0004} & 0.005 & 3: Peak \\
9{,}000 & 13{,}793:1 & 0.0004 & 0.005 & 4: BD stabilization \\
9{,}500 & 12{,}846:1 & 0.0004 & 0.005 & 4: Oscillation \\
\textbf{10{,}000} & \textbf{13{,}456:1} & \textbf{0.0005} & 0.005 & \textbf{4: Final} \\
\bottomrule
\end{tabular}
\caption{CW-001 trajectory: co/cross ratio, Fiedler eigenvalue, and
contrastive weight across four dynamical phases.}
\label{tab:cw001-trajectory}
\end{table}

\textbf{Phase~1 (steps 0--1{,}500): Rapid initial growth.}
The Steersman decays $c_w$ from 0.020 to the floor of 0.005 within
1{,}500~steps as the co/cross ratio rises rapidly. BD structure forms
quickly during this high-$c_w$ window, reaching ${\sim}1{,}800$:1.

\textbf{Phase~2 (steps 1{,}500--4{,}200): Apparent plateau.}
With $c_w$ floored at 0.005, the co/cross ratio oscillates between
1{,}500:1 and 2{,}000:1 for nearly 3{,}000~steps. The Fiedler
eigenvalue rebounds from its Phase~1 minimum, rising to 0.0020---indicating
that spectral connectivity is \emph{increasing} even as BD strength
appears stalled. This phase was initially interpreted as a terminal
state (adaptive decay defeating whisper-strength). It is not.

\textbf{Phase~3 (steps 4{,}600--8{,}300): Exponential breakout.}
The co/cross ratio breaks above 2{,}000:1 at step~4{,}600 and enters an
exponential growth phase: $2{,}252 \to 3{,}860 \to 7{,}884 \to 12{,}921
\to 15{,}183$:1 (peak at step~8{,}300). The Fiedler eigenvalue declines
monotonically from 0.0019 to 0.0004, entering the same regime as the
aggressive-$c_w$ experiments. The growth rate \emph{accelerates} with
accumulated BD strength, suggesting a positive feedback mechanism: once
block-diagonal structure exceeds a critical threshold (${\sim}2{,}000$:1),
the contrastive loss becomes more effective at reinforcing the existing
partition, creating a self-amplifying cycle.

\textbf{Phase~4 (steps 8{,}300--10{,}000): BD regime stabilization.}
The co/cross ratio oscillates between 12{,}145:1 and 15{,}183:1
(mean~13{,}288:1) without a clear growth trend. The Fiedler eigenvalue
stabilizes at 0.00042--0.00048. Validation loss continues declining
(0.4034 $\to$ 0.4017), indicating that the \emph{language model} is still
improving even as the \emph{bridge structure} has reached equilibrium.
This decoupling---structural stasis while perplexity improves---suggests
the bridge has settled into an attractor from which the remaining
training budget refines weights within the established topology rather
than reshaping it.

\subsubsection{Speed, Destination, or Both?}

At step~10{,}000, CW-001's Fiedler eigenvalue (0.00046) is within one
order of magnitude of H-ch6's final Fiedler (0.00009). Both are deep
in the BD regime. But the co/cross ratios tell a more nuanced story:

\begin{table}[h]
\centering
\small
\begin{tabular}{@{}lrrr@{}}
\toprule
Metric & H-ch6 ($c_w = 0.1$ adaptive) & CW-001 ($c_w = 0.005$ floor) & Ratio \\
\midrule
Co/cross & 70{,}404:1 & 13{,}456:1 & $5.2\times$ \\
Fiedler & 0.00009 & 0.00046 & $5.1\times$ \\
Val loss & ${\sim}0.40$ & 0.4017 & ${\sim}$same \\
\bottomrule
\end{tabular}
\caption{CW-001 vs.\ H-ch6 at matched step count (10{,}000 steps).}
\label{tab:cw001-comparison}
\end{table}

A $5.2\times$ gap persists at matched step count. CW-001's Phase~4
stabilization---oscillating at 12{,}145--15{,}183:1 for 2{,}000~steps
with no growth trend---raises the question: is this a transient plateau
(like Phase~2, which lasted 2{,}700~steps before breaking out) or a
genuine attractor at this $c_w$ floor?

Two interpretations compete. Under the \textbf{speed hypothesis}, CW-001
simply needs more training steps; the Phase~$2 \to 3$ breakout precedent
shows the system can appear stalled and then resume exponential growth.
Under the \textbf{attractor hypothesis}, the contrastive floor sets not
only the rate but the ceiling of structural formation. The $5.2\times$
co/cross gap and the $5.1\times$ Fiedler gap are proportional, consistent
with a $c_w$-dependent attractor strength.

The 24C-002 experiment (\S\ref{sec:24cell}) resolves this question in
favour of the attractor hypothesis. At $n{=}24$ with adaptive $c_w$
settling at $0.025$, the final co/cross is $5{,}983$:1---$6\times$ below
the fixed-$c_w{=}0.1$ run (24C-001, $35{,}808$:1) at matched step count.
The contrastive weight governs not only speed but the \emph{depth} of BD
formation. Higher $c_w$ produces a deeper attractor.

Both interpretations share the same novel core: \emph{any nonzero $c_w$
produces BD structure}, the breakout mechanism at ${\sim}2{,}000$:1 is a
genuine phase transition, and the three-phase incubation trajectory
has no prior art in the adapter literature.

\subsubsection{The Phase~2 Plateau as Incubation}

The 3{,}000-step plateau (Phase~2) is not a failure state. During this
period, the Fiedler eigenvalue rises from 0.0004 to 0.0020---the bridge
is building spectral connectivity while block-diagonal strength appears
stalled. We hypothesize that Phase~2 represents an \emph{incubation
period} in which the bridge accumulates the spectral prerequisites for
exponential BD growth. The Fiedler rebound creates a connected substrate;
Phase~3's breakout then partitions this substrate into blocks.

This two-stage process (connect, then partition) mirrors the
spectral-attractor-then-bifurcation sequence observed across the full
experimental programme (\S\ref{sec:spectral-attractor}--\ref{sec:negative-controls}):
spectral-only training creates connectivity (the attractor); contrastive
loss then breaks this connectivity into directed blocks. CW-001
reproduces this sequence \emph{within a single run} at low~$c_w$.

\subsubsection{Implications}

\begin{enumerate}
\item \textbf{No minimum $c_w$ threshold for BD.} Any nonzero $c_w$
  eventually produces BD structure---the FI-004 ``cliff'' at $c_w = 0$
  is real, but $c_w = 0.005$ is far above it.
\item \textbf{Adaptive decay is not a bottleneck.} The Steersman's
  adaptive Control Law~2, which reduces $c_w$ when BD is detected, does
  not prevent BD formation---it merely slows it. The floor mechanism
  ensures a minimum contrastive signal is always present.
\item \textbf{Training budget trades against $c_w$.} A practitioner with
  abundant compute can use lower $c_w$ for potentially better spectral
  properties; a practitioner needing fast convergence can use higher
  $c_w$. Whether the slow regime produces qualitatively better adapters
  (lower Fiedler $=$ more connected $=$ better information flow) is an
  open empirical question.
\item \textbf{The breakout threshold (${\sim}2{,}000$:1) is a phase
  transition.} The co/cross ratio of ${\sim}2{,}000$:1 marks the onset
  of self-amplifying BD growth at $c_w = 0.005$. Pending verification
  across different $n$ and topologies, this may be universal.
\item \textbf{BD regime stabilization decouples from language modeling.}
  Validation loss continues declining (0.4034 $\to$ 0.4017) during
  Phase~4 while co/cross oscillates without growth. The bridge structure
  reaches equilibrium before the language model does---the final training
  budget refines representations within established topology.
\end{enumerate}


\subsection{Limitations}
\label{sec:disc-limitations}

Several limitations qualify the present results.

\textbf{Pair-based specifications only.}
All topology specifications tested partition channels into co-axial
pairs. More complex topologies---triplets, cycles, cliques of
size~${>}\,2$, hierarchical nesting---would require extending the
contrastive loss formulation. Whether the Steersman can program
non-pair-based structures is untested.

\textbf{Single task family.}
All experiments use instruction-following fine-tuning on
alpaca-cleaned. Whether the four-regime taxonomy, spectral attractor,
and geometric selectivity transfer to other task types
(classification, summarization, code generation) is unknown.

\textbf{Single model family.}
Topology programming experiments use TinyLlama~1.1B. Paper~3
established scale consistency for RD contrastive at 1.1B and 7B, but
the generalization of 4D programming, geometric selectivity, and the
full four-regime taxonomy to larger models is untested.

\textbf{No formal optimality proof.}
We demonstrate that the Steersman programs a given topology, but we
do not prove that any particular topology is optimal for a downstream
task. The relationship between topology specification and task
performance remains to be characterized.

\textbf{24-cell CL2 logging blind spot.}
The 24-cell experiment (24C-001) completed 10{,}000 steps. A logging bug
in Control Law~2 (co/cross tracking gated to $n{=}6$) silenced real-time
co/cross after step~4{,}300 ($5{,}673$:1). Post-hoc recovery (PC-001)
from all 100 saved bridge checkpoints revealed that co/cross grew to
$35{,}808$:1---6.3$\times$ beyond the last real-time measurement. The
contrastive loss itself was applied correctly throughout (via the
cybernetic training script, which has no $n$-gate); only the metric
logging was affected. The bug has been documented for correction.

\textbf{FI-003 sign persistence preliminary.}
The sign persistence experiment (\S\ref{sec:disc-homeostasis}) has
1{,}200~steps of LM-only data (12~checkpoints). The two-timescale decay
trajectory is clear, but the asymptotic equilibrium co/cross ratio
(expected to converge to ${\sim}1$:1) and the long-term validation
loss impact are not yet determined.

\textbf{No theoretical account of the attractor or selectivity.}
The spectral attractor at Fiedler $\approx 0.09$ is empirically
consistent but lacks a theoretical derivation. Similarly, we observe
that only polytope-derived pair specifications produce block-diagonal
structure, but provide no formal characterization of which
specifications constitute valid geometric programs. Both are open
questions for future theoretical work, possibly through random matrix
theory (attractor) or group-theoretic analysis of the contrastive loss
landscape (selectivity).
